\subsection{Comentarios de Línea y de Bloque}

\subsubsection{Regla: preferir comentarios de línea breves}

\textbf{Regla:} Se usa un comentario de línea con \texttt{//} para una explicación breve y excepcional. Se evita el comentario de bloque \texttt{/* ... */}; si la explicación requiere varias líneas, se mejora el diseño o se traslada la explicación al documento técnico correspondiente. Ningún comentario conserva código deshabilitado (Microsoft, 2025e). El comentario de documentación \texttt{///} no es un comentario de bloque y se rige por la Sección 7.

\textbf{Código correcto:}
\begin{lstlisting}[style=csharpstyle]
// Saves from version 1.x store the board without the wall count.
int wallCount = saveData.WallCount ?? DefaultWallCount;
\end{lstlisting}

\textbf{Código incorrecto:}
\begin{lstlisting}[style=csharpstyle]
/* Read the wall count from the save data.
   If it is missing, use the default value.
   Then assign it to the local variable. */
int wallCount = saveData.WallCount ?? DefaultWallCount;
// int wallCount = LegacyReader.ReadWalls(saveData);
\end{lstlisting}

\subsubsection{Regla: comentarios en scripts SQL}

\textbf{Regla:} Un script SQL admite un único comentario de bloque como encabezado, que describe el propósito del script. Un comentario de línea \texttt{--} junto a una columna se admite solo cuando documenta una regla que el nombre no expresa: un rango de valores, el origen del dato o el requisito que la exige. Prohibido el comentario que repite el nombre de la columna.

\textbf{Código correcto:}
\begin{lstlisting}[language=SQL,basicstyle=\ttfamily\small,frame=single,breaklines=true,columns=fullflexible]
PlayerCount     TINYINT  NOT NULL,  -- 2 or 4, never 3 (CU-17 RN-02).
WallsPerPlayer  TINYINT  NOT NULL,
\end{lstlisting}

\textbf{Código incorrecto:}
\begin{lstlisting}[language=SQL,basicstyle=\ttfamily\small,frame=single,breaklines=true,columns=fullflexible]
GameModeId      TINYINT  NOT NULL,  -- Mode identifier.
WallsPerPlayer  TINYINT  NOT NULL,  -- Walls per player.
\end{lstlisting}
